% ==========================================================================
% GENERATED by scripts/split_paper.py. Do not edit this file.
%
% Source: imported_papers/sds_hounsfield_features/main.tex
% Part 5: Results
%
% Edits belong upstream, in the published paper. To change what the
% chapter includes, change imports/hounsfield.toml and re-run
% the splitter.
% ==========================================================================
\section{Results}

\subsection{Reduced data needs}

\begin{figure}[t]
    \centering
    \includegraphics[width=1.0\linewidth]{figures/results/less_data_perspective_ji.png}
    \caption{Our model using HU ranges as additional features shows consistent improvement \emph{across different dataset sizes} over the baseline without them. The horizontal blue line indicates the \emph{U-Net base, HU} baseline performance using the entire dataset.
    The vertical orange line indicates that our approach surpasses this performance using \formatpercentage[round-precision=0]{0.50} of the dataset.
    % set is necessary to match the performance of the ``\emph{U-Net base, HU}'' baseline trained on all data.
    }
    \label{fig:robustness-data-size}
\end{figure}

\Cref{fig:robustness-data-size} presents the Dice score for the U-Net-base model across varying dataset fractions, comparing performance with and without the proposed HU ranges as additional features.
The results show our approach consistently improves performance by approximately \SIrange{1}{2}{\percent} across dataset fractions except for the \SI{10}{\percent} fraction.
This smallest dataset fraction only contains $4$ patients, which increases variability in cross-validation and results in wider confidence intervals.

More importantly, as highlighted in \cref{fig:robustness-data-size}, the model utilizing HU ranges surpasses the performance of the baseline when only using \SI{50}{\percent} of the training data. This finding underscores the \emph{data efficiency} of our approach, suggesting that incorporating HU ranges not only enhances segmentation accuracy but also significantly reduces the amount of labeled data needed to achieve a given performance level. Consequently, our method offers substantial practical benefits by mitigating annotation costs and improving data efficiency in training deep learning models in and for clinical practice.

\subsection{Performance improvements}

%
\begin{table*}[t]
    \caption{Experimental results. Score ranges represent the \SI{95}{\percent} confidence interval. Dice score, Precision, Recall, and F$_1$-Score are all macro averages across tissue classes.}
    \centering
    \begin{tabular}{lrrrrr}
        \toprule
        \multicolumn{1}{c}{\textbf{Model}} 
        & \multicolumn{1}{c}{\textbf{Number of parameters}} 
        % & \multicolumn{1}{c}{\textbf{Train Time Factor}} 
        & \multicolumn{1}{c}{\textbf{Dice score}} 
        & \multicolumn{1}{c}{\textbf{Precision}} 
        & \multicolumn{1}{c}{\textbf{Recall}} 
        & \multicolumn{1}{c}{\textbf{F$_1$-Score}} \\
        \midrule
        U-Net-small, HU                                   
        & \num{1943829} 
        % & \num{1.00}            
        % & \formatpercentage{0.8580}\ \ $\pm$\formatpercentage{0.0152}
        &  \hspace{15pt} \formatpercentage{0.9303}\ \ $\pm$\formatpercentage{0.002597}
        &  \hspace{5pt} \formatpercentage{0.9246}\ \ $\pm$\formatpercentage{0.01549}
        &  \hspace{5pt} \formatpercentage{0.9272}\ \ $\pm$\formatpercentage{0.009421} 
        &  \hspace{5pt} \formatpercentage{0.9252}\ \ $\pm$\formatpercentage{0.003493} \\
        U-Net-small, HU + HU ranges                       
        & \num{1944549} 
        % & \num{1.17}               
        % & \formatpercentage{0.8950}\ \ $\pm$\formatpercentage{0.0081} 
        & \formatpercentagebold{0.9467}\ \ $\pm$\formatpercentage{0.002802} 
        & \formatpercentagebold{0.9357}\ \ $\pm$\formatpercentage{0.003558} 
        & \formatpercentagebold{0.953}\ \ $\pm$\formatpercentage{0.004187} 
        & \formatpercentagebold{0.944}\ \ $\pm$\formatpercentage{0.002709} \\
        \midrule
        U-Net-base,\ \ HU                                    
        & \num{7765541}
        % & \num{1.00}  % \num{2.01}              
        % & \formatpercentage{0.8681}\ \ $\pm$\formatpercentage{0.0109} 
        & \formatpercentage{0.9359}\ \ $\pm$\formatpercentage{0.003291} 
        & \formatpercentage{0.9266}\ \ $\pm$\formatpercentage{0.004032} 
        & \formatpercentage{0.9378}\ \ $\pm$\formatpercentage{0.007159} 
        & \formatpercentage{0.9319}\ \ $\pm$\formatpercentage{0.003644} \\
        U-Net-base,\ \ HU + HU ranges                        
        & \num{7766981} 
        % & \num{1.08} % \num{2.17}   
        % & \formatpercentage{0.8980}\ \ $\pm$\formatpercentage{0.0061} 
        & \formatpercentagebold{0.9481}\ \ $\pm$\formatpercentage{0.002738} 
        & \formatpercentagebold{0.9403}\ \ $\pm$\formatpercentage{0.003056}
        & \formatpercentagebold{0.95}\ \ $\pm$\formatpercentage{0.008919}
        & \formatpercentagebold{0.9449}\ \ $\pm$\formatpercentage{0.002831}  \\
        \midrule
        U-Net-large, HU                                   
        & \num{31042629} 
        % & \num{1.00} % \num{5.82}          
        % & \formatpercentage{0.8655}\ \ $\pm$\formatpercentage{0.0087} 
        & \formatpercentage{0.9312}\ \ $\pm$\formatpercentage{0.005912} 
        & \formatpercentage{0.9263}\ \ $\pm$\formatpercentage{0.01057} 
        & \formatpercentage{0.9273}\ \ $\pm$\formatpercentage{0.02222}
        & \formatpercentage{0.9261}\ \ $\pm$\formatpercentage{0.007489} \\
        U-Net-large, HU + HU ranges                       
        & \num{31045509}          
        % & \num{1.02} % \num{5.93}               
        % & \formatpercentage{0.8913}\ \ $\pm$\formatpercentage{0.0039} 
        & \formatpercentagebold{0.9455}\ \ $\pm$\formatpercentage{0.001213} 
        & \formatpercentagebold{0.9316}\ \ $\pm$\formatpercentage{0.003021} 
        & \formatpercentagebold{0.953}\ \ $\pm$\formatpercentage{0.006178} 
        & \formatpercentagebold{0.9419}\ \ $\pm$\formatpercentage{0.001499}  \\
        \bottomrule
    \end{tabular}
    \label{tab:results} \label{tbl:results}
\end{table*}

\Cref{tbl:results} summarizes the experimental results, comparing the performance of U-Net models of varying sizes with and without our proposed approach of incorporating HU ranges as additional features to learn from.
Across all model sizes, the inclusion of HU ranges consistently improves segmentation performance across multiple metrics, with the Dice score improving by approximately \SIrange{1}{2}{\percent}, reducing the remaining error by \SI{20}{\percent}. 
The improvements hold consistently across different U-Net model sizes and dataset fractions, demonstrating the applicability across diverse training settings.

Our baseline results without our proposed additional inputs align with previously reported findings on similar datasets.
For example, \cite{ha_development_2021} achieve a Dice score of \formatpercentage[round-precision=0]{0.97} using their internal dataset of $1496$ CT scans of the L3 vertebra from $922$ patients, segmenting the tissue types SM, SAT, and VAT.
Our \emph{U-Net-base, HU} model achieves a comparable Dice score of \formatpercentage[round-precision=0]{0.9492} on the same tissue types (i.e., when ignoring IMAT) despite using a smaller dataset of only $176$ CT scans from $44$ patients for training and evaluation.

By expanding the number of input channels with engineered masks, the first convolution has few additional parameters, leading to increased computational demand and memory needs.
Empirically, we found that the training time was at most $1.17\times$ higher than the network without the HU ranges.
This efficiency makes our approach both computationally feasible and practically advantageous.  

\subsection{Performance per tissue type}

\begin{figure}[t]
\centering
\includegraphics[width=1.0\linewidth]{figures/results/base_tissue_types.png}
\caption{Dice score \emph{per tissue type} and background (BG). Giving the model a physiology-informed inductive bias via HU ranges-based masks keeps (BG, SAT) or improves performance, with IMAT improvements being of high clinical value.}
\label{fig:tissue-types}
\end{figure}

\Cref{fig:tissue-types} shows the dice score per tissue type for the U-Net-base model, comparing performance with and without the inclusion of HU ranges.
While the segmentation performance remains nearly identical for background and SAT, a modest improvement is observed for SM and VAT.
However, a substantial performance increase of approximately \SI{5}{\percent} is achieved for IMAT.


IMAT is recognized as an important biomarker for various medical diagnoses \cite{lim_inter-muscular_2019, suka_aryana_important_2023}, yet it remains challenging to segment it accurately \cite{hemke_deep_2020, paris_automated_2020} because its regions are smaller and more spread out than other tissues.
Thus, the incorporation of HU ranges not only enhances overall segmentation performance but also significantly improves the detection of IMAT, a clinically relevant yet difficult-to-identify tissue type.
